\documentclass[11pt,a4paper]{scrartcl} %https://ctan.org/pkg/scrartcl : script-article

\usepackage[T1]{fontenc} %https://ctan.org/pkg/fontenc : encodings manager
\usepackage[ngerman]{babel} %https://ctan.org/pkg/ : Sprachenunterstützung
\usepackage{microtype} %https://ctan.org/pkg/ : Erweiterte Unterstützung für Scrtiften
\usepackage[hidelinks]{hyperref} %https://ctan.org/pkg/ : cross-referencing in LaTeX 


% Gliederung bis subparagraph nummerieren
\setcounter{secnumdepth}{5}

% Gliederung bis subparagraph im
% Inhaltsverzeichnis anzeigen
\setcounter{tocdepth}{5}


\begin{document}

\begin{titlepage}
  \centering
  \vspace*{3cm} % vertical space of 3 cm
  {\Huge\bfseries Meine Forschungsarbeit\par} %\Huge => huge font; \bfseries => "bold" i.e. "fett"
  \vspace{1.5cm}
  {\Large Untertitel oder Forschungsfrage\par} %\Large => karge font, Large<Huge
  \vfill %fügt einen unendlich dehnbaren vertikalen Abstand ein, der den verbleibenden leeren Platz auf einer Seite komplett ausfüllt
%  \vspace{12.5cm}
  Vorname Nachname\\
  Schule / Arbeitsgemeinschaft\\
  \today %Kompilierungsdatum
\end{titlepage}

\begin{abstract} %
  In zwei bis vier Sätzen: Frage, Methode, Ergebnis.
\end{abstract}

\tableofcontents
\clearpage %Beginnt eine neue Seite ähnlich wie \newpage

\section{Einleitung}
Hier steht die Forschungsfrage.
\section{Methode}
\section{Experimentelle Untersuchung}
\subsection{Versuchsaufbau}
\subsubsection{Messinstrumente}
\paragraph{Temperatursensoren}
Beschreibung der Temperatursensoren.
\subparagraph{Kalibrierung}
Beschreibung der Kalibrierung.


Hier wird der Versuchsaufbau beschrieben.
\section{Ergebnisse}
Hier stehen Messwerte und Beobachtungen.
\section{Diskussion und Fazit}
Was bedeuten die Ergebnisse?
\section*{Verschiedenes} % '*' nach \section, \subsection usw. schließen die Numerierung für diesen Abschnitt ab 
ggf. zusätzliche relevante Information
%\addsec{Vorwort}
%Hier steht der Text des Vorworts. Es hat keine Nummer, wird aber im 
%Inhaltsverzeichnis korrekt aufgeführt und der Hyperlink springt an die richtige Stelle.

%TODO: \clearpage vs \newpage mit Beispiel
%TODO: \addsec{Vorwort}
%TODO: Problem: Formel zu lang!!! was tun
%TODO: Formeln mit backslash? was ist besser?
%TODO: Chemische Abbildungen - geskriptet


\end{document}
